% Organização sugerida para uma introdução de relatório de estágio

\chapter{Introdução}\label{sec:1}

% \textbf{(MÍNIMO UMA PÁGINA)}

% A introdução do trabalho deverá contextualizar a proposta de desenvolvimento do estágio. 

% Contextualizar deve ser entendido aqui como uma explicação de como o acadêmico chegou à esta proposta de estágio supervisionado.

% Lembre que na introdução de seu trabalho a seguinte pergunta deverá ser  respondida ao leitor : “se eu continuar a ler o trabalho daqui para frente o que vou encontrar no seu texto?”

% Exemplo:

% (resumido e fictício apenas para fins de exemplo)

% “A Lei Complementar no 101, de 4 de maio de 2000, intitulada Lei de Responsabilidade Fiscal - LRF,
% estabelece normas de finanças públicas voltadas para a responsabilidade na gestão
% fiscal, mediante ações em que se previnam riscos e corrijam desvios capazes de afetar o
% equilíbrio das contas públicas, destacando-se o planejamento, o controle, a
% transparência e a responsabilização como premissas básicas.” Tesouro Nacional, 2008

% A partir da publicação da LRF ficou evidente que a utilização de recursos
% tecnológicos para acompanhamento da gestão se faz necessária nas diversas
% esferas do poder público.

% ...

% O presente projeto de estágio propõe a análise e o desenvolvimento de uma
% aplicação de controle orçamentário para a Prefeitura Municipal do Vale dos Sonhos.

% O projeto de estágio descreve os objetivos do trabalho, detalha a proposta do
% sistema a ser desenvolvido, faz um breve resgate da teoria que ampara o sistema
% proposto e apresenta o cronograma e ferramental de hardware e software estimado
% para o desenvolvimento do trabalho.

% \section{PROPOSTA E OBJETIVOS}

%Quais objetivos (geralmente 1 objetivo geral e no mínimo 5 específicos) o acadêmico pretende atingir com o desenvolvimento do estágio.

% Evite aqui objetivos como: “o objetivo deste trabalho é concluir o curso de
% graduação”, o objetivo a ser abordado neste tópico é o objetivo do sistema proposto.


% Exemplo (resumido e fictício apenas para fins de exemplo)

%    O presente projeto tem como objetivo principal (ou objetivo geral) propor uma
% solução de acompanhamento do orçamento do município de Vale dos Sonhos.
% Como objetivos específicos o sistema irá proporcionar a o cadastramento e a consulta do orçamento municipal aprovado pela Câmara de Vereadores de Vale dos
% Sonhos, tão bem quanto o acompanhamento on-line dos empenhos e pagamentos
% efetuados

% (os objetivos gerais e específicos poderão ser feitos na forma de tópicos se melhor convier ao aluno)

% \section{JUSTIFICATIVA}

% Qual é a proposta do seu sistema? Como (e porque) ela se justifica como uma proposta válida?

% Neste item você dever fazer uma breve descrição do seu sistema (proposta) e porque ele é importante para a empresa cedente (justificativa).

% Exemplo :

% (este exemplo está resumido e é fictício)

% O sistema proposto permitirá o controle dos gastos públicos por meio da inclusão dos pagamentos efetuados pelo setor financeiro da Prefeitura do Vale dos Sonhos disponibilizando o acesso ao orçamento municipal imediatamente aos contribuintes que efetuarem consultas na Internet. (PROPOSTA)

% “A LRF cria condições para a implantação de uma nova cultura gerencial na gestão dos recursos públicos e incentiva o exercício pleno da cidadania, especialmente no que se refere à participação do contribuinte no processo de acompanhamento da aplicação dos recursos públicos e de avaliação dos seus resultados”, desta forma o presente projeto se justifica por atender a uma demanda da Prefeitura Municipal de Vale dos Sonhos buscando a transparência na gestão dos recursos públicos do município. (JUSTIFICATIVA)


Uma instituição de ensino superior enfrenta diversos desafios administrativos para poder garantir uma boa qualidade no ensino oferecido. O avanço da tecnologia da informação vem proporcionando diversas melhorias para a otimização de processos administrativos, com redução de falhas, riscos e custos operacionais.

Em função do tempo o gerenciamento de salas de uma instituição de ensino pode acabar se tornando uma tarefa complexa, devido às necessidades como as trocas de horários, alta demanda e controles de restrições. Assim, processos realizados de forma manual como a utilização de planilhas podem acabar resultando em conflitos de agendamento, retrabalho e dificuldades no acompanhamento das reservas realizadas.

O presente relatório de estágio descreve o desenvolvimento realizado de uma aplicação de controle de agendamentos de salas para o Núcleo de Prática de Informática - NPI.

Este relatório descreve os objetivos alcançados, detalha os casos de uso implementados, os recursos e o ambiente de desenvolvimento utilizados, bem como os resultados obtidos ao longo do estágio.

\section{PROPOSTA E OBJETIVOS}

% O presente projeto tem como objetivo principal (ou objetivo geral) propor uma solucao de agendamento de salas para a unifil. o Sistema irá proporcionar controle de resolução de conflito e cadastros de salas e usuários, controle de históricos de alterações, uma agenda centralizada da locação dos recursos e geração de relatórios.

O presente trabalho teve como objetivo principal (ou objetivo geral) desenvolver uma solução de agendamento de salas para a UniFil. Os seus objetivos específicos foram:

um fluxo padrão para agendamento.

controle de históricos de alterações

\section{JUSTIFICATIVA}

O sistema em desenvolvimento permite o controle de alocação de salas por meio de uma centralização das informações e automatização do processo de agendamento. O sistema está substituindo o processo anterior realizado em planilhas, garantindo maior confiabilidade e eficiência operacional.